\section{Producción audiovisual en directo}

La producción audiovisual en directo engloba el conjunto de procesos técnicos y
organizativos que permiten capturar, mezclar y distribuir contenido de vídeo y audio
en tiempo real, sin posibilidad de edición posterior antes de su emisión. A diferencia
de la producción diferida, donde los errores se corrigen en postproducción, el directo
exige que todos los elementos del sistema funcionen de forma sincronizada y sin
interrupciones durante toda la duración del evento.

\subsection{Flujo de trabajo de una realización}

El flujo de trabajo típico de una realización en directo sigue una cadena bien definida:
las fuentes de vídeo (cámaras, ordenadores de presentación, reproductores de vídeo)
generan señales que llegan al \textit{mezclador} o \textit{switcher}, donde el
realizador selecciona en cada instante qué fuente se emite al programa. La señal de
programa resultante se envía simultáneamente al sistema de grabación y al sistema de
distribución (streaming, emisión broadcast o proyección en sala).

Sobre la señal de programa pueden superponerse \textit{overlays} gráficos: logotipos,
barras de información, rótulos con el nombre del ponente (\textit{lower thirds}) o
señales indicadoras del estado de la emisión. Esta superposición puede realizarse
antes o después de la mezcla, dependiendo de la arquitectura del sistema.

\subsection{Terminología fundamental}

Los conceptos clave en producción en directo que aparecen a lo largo de este trabajo son:

\begin{itemize}
  \item \textbf{Switcher / mezclador}: dispositivo o software que selecciona y combina
  las fuentes de vídeo para producir la señal de programa.
  \item \textbf{Corte} (\textit{cut}): transición instantánea de una fuente a otra,
  sin efecto intermedio.
  \item \textbf{Composite}: modo en que dos o más fuentes se combinan visualmente en
  pantalla simultáneamente (p. ej., \textit{picture-in-picture}, pantalla dividida).
  \item \textbf{Overlay}: elemento gráfico superpuesto sobre la señal de vídeo sin
  reemplazarla, como logotipos o rótulos.
  \item \textbf{Stream blanker}: mecanismo que sustituye la señal de programa por una
  señal alternativa (pausa o corte de señal) cuando la emisión debe interrumpirse
  temporalmente sin detener el sistema.
  \item \textbf{Preview} / \textbf{Program}: el bus de programa contiene la señal que
  se emite al exterior; el bus de preview muestra la fuente que se emitirá a
  continuación, permitiendo al realizador prepararla.
\end{itemize}

\subsection{Producción broadcast tradicional frente a producción IP}

Históricamente, la producción en directo se ha basado en infraestructura hardware
dedicada: mezcladores de vídeo profesionales (Ross, Blackmagic ATEM, Grass Valley),
matrices de enrutamiento y cableado SDI (\textit{Serial Digital Interface}).
Este modelo ofrece latencias subimagen y alta fiabilidad, pero implica un coste
elevado y una infraestructura poco portable.

La producción IP o basada en software reemplaza parte de ese hardware por procesamiento
en CPU/GPU de propósito general. Herramientas como Voctomix, OBS Studio o vMix ejecutan
el mezclado íntegramente por software, lo que permite desplegar un estudio de producción
completo sobre hardware convencional y reducir los costes de infraestructura de forma
significativa. La contrapartida es una mayor dependencia del sistema operativo y una
latencia algo superior, que en la mayoría de casos de uso ---retransmisión de
conferencias, eventos académicos, producciones comunitarias--- resulta perfectamente
aceptable.
